\section{相对论}

\subsection{洛伦兹变换}

\subsubsection{狭义相对论基本原理}

爱因斯坦提出两个基本假设：

\begin{enumerate}
    \item 狭义相对性原理：物理定律在一切惯性系中都具有相同形式；
    \item 光速不变原理：在一切惯性系中，真空中光速 $c$ 都是相同的。
\end{enumerate}

\subsubsection{洛伦兹变换}

对于两个参考系 $S, S'$，若 $S'$ 以速度 $v$ 沿 $x$ 轴正方向运动，定义时间零点为两坐标系的原点重合的时刻。则洛伦兹正变换的变换式为：
$$
\begin{dcases}
    x' = \gamma (x - vt) \\
    y' = y \\
    z' = z \\
    t' = \gamma \ab(t - \frac{v x}{c^2}) \\
\end{dcases}
$$
其中 $\gamma = \frac{1}{\sqrt{1 - v^2 / c^2}}$，$c$ 为真空光速。

\subsection{相对论时空观}

\subsubsection{同时性的相对性}

在某个参考系下同时发生的事件，在另一个参考系不一定同时发生。两个事件之间的时间关系同样不一定保持。

\subsubsection{时序与因果关系}

对于两个具有因果关系的时间，其时序不会颠倒。因为两事件之间信息的传递速度不超过光速，即 $\frac{\Delta x}{\Delta t} \le c$。因此
$$
\Delta t' = \gamma \ab(\Delta t - \frac{v \Delta x}{c^2}) = \gamma \Delta t \ab(1 - \frac{v}{c} \frac{\Delta x}{\Delta t}) \ge 0
$$

\subsubsection{尺缩效应}

在惯性系中测量一个物体的长度，其在运动方向上的长度（相对于静止时）长度缩短。

这里测量运动物体必须同时测量两端的坐标。
